% ============================================================
% Bölüm 8: Bizim Problemimiz -- 3D Mixed Convection
% ============================================================

\chapter{Bizim Problemimiz: 3D Mixed Convection}
\label{ch:bizim_problem}

\"Onceki bölümlerde teorik temelleri kurduk: Navier--Stokes denklemleri, Boussinesq yaklaşımı, spektral yöntemler, türbülans ve LES, neural operator kavramları. Şimdi bu bilgilerin tamamını somut bir probleme uyguluyoruz.

Bu bölüm, projenin fiziksel çerçevesini eksiksiz tanımlar: geometri, sınır koşulları, zorlama mekanizması, sıcaklık ayrıştırması, boyutsuz denklemler, parametre uzayı, CFL analizi ve başlangıç koşulları.


% ============================================================
\section{Problem Tanımı ve Geometri}
\label{sec:problem_tanim}

\subsection{Domain}

Çözüm alanı, üç boyutlu dikdörtgen bir kutudur:
\begin{equation}
    \Omega = [0, L_x] \times [0, L_y] \times [0, L_z], \quad
    L_x = 6, \quad L_y = 10, \quad L_z = 4.
    \label{eq:domain}
\end{equation}
Burada $y$-ekseni dikey (yerçekimi yönü), $x$-ekseni akışın birincil yönü ve $z$-ekseni enine (spanwise) yöndür. Tüm uzunluklar boyutsuzdır; referans uzunluk $H$ alınmıştır.

Domain'in en-boy oranları $L_x : L_y : L_z = 3:5:2$ olarak seçilmiştir. Dikey boyutun ($L_y = 10$) büyük seçilmesinin nedeni: kaldırma kuvveti kaynaklı büyük ölçekli konveksiyon hücrelerinin (plume) yeterli alan bulmasıdır. Yatay boyutlar ($L_x = 6$, $L_z = 4$) ise Kolmogorov forcing'in yeterli sayıda dalga boyunu barındırması ve spanwise yapıların gelişebilmesi için uygundur.

\subsection{Sınır Koşulları}

\textbf{Tüm yönlerde periyodik:}
\begin{equation}
    \vect{u}(0, y, z, t) = \vect{u}(L_x, y, z, t), \quad
    \vect{u}(x, y, 0, t) = \vect{u}(x, y, L_z, t), \quad
    \vect{u}(x, 0, z, t) = \vect{u}(x, L_y, z, t).
    \label{eq:periodic_bc}
\end{equation}
Bu tamamen periyodik yapı, FFT tabanlı spektral yöntemlerin doğrudan uygulanmasına olanak tanır: türevler, Poisson çözücü ve dealiasing işlemleri Fourier uzayında analitik olarak yapılır.

\textbf{Sıcaklık koşulları:}
\begin{itemize}
    \item Alt yüzey ($y = 0$): sıcak, $T = T_{\text{hot}} = 20\,^{\circ}\mathrm{C}$
    \item Üst yüzey ($y = L_y$): soğuk, $T = T_{\text{cold}} = 0\,^{\circ}\mathrm{C}$
    \item Sıcaklık farkı: $\Delta T = T_{\text{hot}} - T_{\text{cold}} = 20\,^{\circ}\mathrm{C}$
\end{itemize}
Bu sıcaklık koşulu, Rayleigh--B\'{e}nard konveksiyonunun sürücüsüdür. Alttaki sıcak akışkan hafifleyerek yükselir, üstteki soğuk akışkan ağırlaşarak alçalır; bu dengesizlik konvektif hücrelere yol açar.

\begin{remark}
Alt sıcak, üst soğuk koşulu ile periyodik sınır koşulu arasında ciddi bir çelişki vardır. Bu çelişki ve çözümü Bölüm~\ref{sec:sicaklik_ayristirma}'da ele alınmaktadır.
\end{remark}

\subsection{Fiziksel Senaryo}

Bu problem, aşağıdaki gerçekçi senaryoları modeller:
\begin{itemize}
    \item \textbf{Dikey atrium:} Çok katlı bir binanın ortasındaki açık hacim. Taban ısıtılmış (yer ısıtma), tavan soğuk; aynı zamanda kapı/pencere açıklıklarından yatay rüzgâr girer.
    \item \textbf{Açık alan:} Güneş'ten ısınan zemin üzerinde, yatay rüzgârın etkisi altında oluşan atmosferik sınır tabaka.
    \item \textbf{Bina aralarındaki rüzgâr koridoru:} Binalar arasında kanallanmış rüzgâr ve zeminle bina yüzeyleri arasındaki sıcaklık farkının örneği.
\end{itemize}

%\begin{figure}[H]
%    \centering
%    \includegraphics[width=0.85\textwidth]{problem_geometry.png}
%    \caption{3D mixed convection domain geometrisi. $L_x \times L_y \times L_z = 6 \times 10 \times 4$ boyutlu periyodik kutu. Alt yüzey sıcak ($T_{\text{hot}}$), üst yüzey soğuk ($T_{\text{cold}}$). Yatay ok, Kolmogorov forcing'in yönünü gösterir.}
%    \label{fig:problem_geometry}
%\end{figure}


% ============================================================
\section{Kolmogorov Forcing}
\label{sec:kolmogorov_forcing}

\subsection{Formülasyon}

Hız alanına uygulanan dış kütle kuvveti (body forcing):
\begin{equation}
    \vect{F}(y) = A \sin\!\left(\frac{k_f \cdot 2\pi y}{L_y}\right) \hat{\vect{e}}_x
    \label{eq:kolmogorov_forcing}
\end{equation}
\begin{itemize}
    \item $A$: genlik (öğrenilebilir parametre, INNATE'te)
    \item $k_f = 1$: zorlama dalga sayısı (tek mod)
    \item $\hat{\vect{e}}_x$: kuvvet $x$-yönündedir
    \item Kuvvet yalnızca $y$'ye bağlıdır: dikey konuma göre farklı şiddette rüzgâr
\end{itemize}

$k_f = 1$ seçimi, domain içinde tek bir sinüs dalgası anlımına gelir: $y = L_y/4$'te kuvvet pozitif ($+x$), $y = 3L_y/4$'te negatif ($-x$). Bu da bir kesme (shear) akışı yaratır.

\subsection{Neden Kolmogorov Forcing?}

Kolmogorov akışı (Kolmogorov flow), türbülans araştırmalarında standart bir benchmark'tir. Seçim nedenleri:

\begin{enumerate}
    \item \textbf{Analitik laminer çözüm:} Düşük $\Rey$ değerlerinde, karılı durum çözümü:
    \begin{equation}
        u_{\text{laminar}}(y) = \frac{A \, \Rey}{(k_f \cdot 2\pi / L_y)^2} \sin\!\left(\frac{k_f \cdot 2\pi y}{L_y}\right)
        \label{eq:kolmogorov_laminar}
    \end{equation}
    Bu analitik çözüm, kodun doğrulanması için vazgeçilmezdir.

    \item \textbf{Periyodik yapıya uyum:} Forcing fonksiyonu doğal olarak periyodiktir, FFT ile tam uyum içindedir.

    \item \textbf{Akademik standart:} Kolmogorov-type forcing, $\Rey$ aşıldığında türbülansa geçişi temiz biçimde üretir. Energy injection rate iyi tanımlıdır: $\varepsilon_{\text{input}} = \langle \vect{u} \cdot \vect{F} \rangle$.

    \item \textbf{Kontrol edilebilir rejim:} $A$ ayarlanarak, zorlanmış konveksiyon şiddeti sürekli değiştirilebilir.
\end{enumerate}

\subsection{Stochastic Mod: Ornstein--Uhlenbeck Süreci}

Gerçek rüzgâr sabit değildir; zamanla dalgalanır. Bu davranışı modellemek için genlik $A(t)$ bir Ornstein--Uhlenbeck (OU) sürecine bağlanır:
\begin{equation}
    dA = -\frac{A - A_0}{\tau_{\text{OU}}} \, dt + \sigma_{\text{OU}} \, dW(t)
    \label{eq:ou_process}
\end{equation}
\begin{itemize}
    \item $A_0$: ortalama genlik (mean reversion hedefi)
    \item $\tau_{\text{OU}}$: korelasyon zamanı (büyük $\to$ yavaş değişim)
    \item $\sigma_{\text{OU}}$: dalgalanma şiddeti
    \item $dW(t)$: Wiener süreci (Brownian motion)
\end{itemize}
OU süreci, sıfıra değil $A_0$'a dönme eğilimindedir; dolayısıyla rüzgâr hiçbir zaman ``ölmez'' ama gerçekçi şekilde dalgalanır. $k_f = 1$ seçimi ile bu mod, büyük ölçekli atmosferik rüzgâr dalgalanmasının basit bir surrogat'ıdır.

%\begin{figure}[H]
%    \centering
%    \includegraphics[width=0.80\textwidth]{kolmogorov_forcing.png}
%    \caption{Kolmogorov forcing profili. Sol: $F_x(y) = A\sin(2\pi y / L_y)$. Sağ: Stochastic modda genlik $A(t)$'nin OU süreci ile zamansal evrimi.}
%    \label{fig:kolmogorov_forcing}
%\end{figure}


% ============================================================
\section{Sıcaklık Ayrıştırması}
\label{sec:sicaklik_ayristirma}

Bu kısım, projenin teknik olarak en kritik noktalarından biridir. Periyodik sınır koşulları ile inhomogen sıcaklık koşulunu bağdaştırmak için sıcaklık alanının \textit{ayrıştırılması} gerekir.

\subsection{Sorun: Periyodik BC ile Sıcaklık Çelişkisi}

Periyodik sınır koşulu, $y$-yönünde şunu zorunlu kılar:
\begin{equation}
    T(x, 0, z, t) = T(x, L_y, z, t).
    \label{eq:periodic_T}
\end{equation}
Ancak fiziksel problemde:
\begin{equation}
    T(x, 0, z, t) = T_{\text{hot}} = 20\,^{\circ}\mathrm{C}, \quad
    T(x, L_y, z, t) = T_{\text{cold}} = 0\,^{\circ}\mathrm{C}.
\end{equation}
Bu iki koşul aynı anda sağlanamaz. Eğer $T$ doğrudan FFT'ye sokulursa, periyodiklik varsayımı ihlal edilir ve Gibbs fenomeni (süreksizlikteki salınım) bütün çözümü bozar.

\subsection{Çözüm: Baz Profil + Pertürbasyon}

Sıcaklık alanı iki parçaya ayrılır:
\begin{equation}
    \boxed{T(x, y, z, t) = T_{\text{base}}(y) + T'(x, y, z, t)}
    \label{eq:T_decomposition}
\end{equation}

\textbf{Baz profil} (sabit, çözüm boyunca değişmez):
\begin{equation}
    T_{\text{base}}(y) = T_{\text{hot}} - \frac{\Delta T}{L_y} \, y
    \label{eq:T_base}
\end{equation}
Bu, alt sıcak -- üst soğuk koşulunu sağlayan lineer profildir. Fiziksel olarak, konveksiyon yokken (tamamen difuzif rejim), sıcaklık alanının yaklaşacağı denge profilidir.

\textbf{Pertürbasyon} (zamanla evrilen, modelin çözdüğü büyüklük):
\begin{equation}
    T'(x, y, z, t) = T(x, y, z, t) - T_{\text{base}}(y)
    \label{eq:T_perturbation}
\end{equation}
$T'$ \textit{periyodiktir}:
\[
    T'(x, 0, z, t) = T(x, 0, z, t) - T_{\text{hot}} = 0, \quad
    T'(x, L_y, z, t) = T(x, L_y, z, t) - T_{\text{cold}} = 0.
\]
Dolayısıyla $T'(x, 0, z, t) = T'(x, L_y, z, t) = 0$ olup periyodiklik koşulu doğal olarak sağlanır. FFT sorunsuz uygulanabilir.

\subsection{Enerji Denklemine Etkisi}

Orijinal enerji denklemi (Boussinesq, boyutsuz):
\begin{equation}
    \pd{T}{t} + (\vect{u} \cdot \gradient) T = \frac{1}{\Rey \, \Pran} \laplacian T.
    \label{eq:energy_original}
\end{equation}
$T = T_{\text{base}}(y) + T'$ yerlerine koyulduğunda:
\begin{align}
    \pd{T'}{t} + (\vect{u} \cdot \gradient)(T_{\text{base}} + T')
    &= \frac{1}{\Rey \, \Pran} \laplacian (T_{\text{base}} + T') \notag \\
    \pd{T'}{t} + (\vect{u} \cdot \gradient) T_{\text{base}} + (\vect{u} \cdot \gradient) T'
    &= \frac{1}{\Rey \, \Pran} \left[ \underbrace{\laplacian T_{\text{base}}}_{=0 \text{ (lineer profil)}} + \laplacian T' \right]. \notag
\end{align}
$T_{\text{base}}(y)$ lineer olduğundan $\laplacian T_{\text{base}} = 0$. Adveksiyon terimi ise:
\begin{equation}
    (\vect{u} \cdot \gradient) T_{\text{base}}
    = u \underbrace{\pd{T_{\text{base}}}{x}}_{=0}
    + v \pd{T_{\text{base}}}{y}
    + w \underbrace{\pd{T_{\text{base}}}{z}}_{=0}
    = v \cdot \left(-\frac{\Delta T}{L_y}\right)
    = -\frac{v \, \Delta T}{L_y}.
\end{equation}
Sonuç olarak, $T'$ için evrim denklemi:
\begin{equation}
    \boxed{
        \pd{T'}{t} + (\vect{u} \cdot \gradient) T'
        = \frac{1}{\Rey \, \Pran} \laplacian T'
        + \frac{v \, \Delta T}{L_y}
    }
    \label{eq:T_prime_equation}
\end{equation}

Sağ taraftaki son terim, $v \cdot \Delta T / L_y$, bir \textbf{kaynak terimdir}. Fiziksel anlamı: dikey hız $v > 0$ (yukarı akış) olduğunda, sıcak akışkan lineer profilden daha sıcak bir bölgeye taşınır $\to$ pozitif $T'$ üretilir. Bu, kaldırma kuvveti ile sıcaklık alanının kenetlenmesini (coupling) sağlar.

\subsection{Pratikte Uygulama}

\begin{enumerate}
    \item \textbf{State değişkeni:} Model yalnızca $T'$'yi tutar ve evrimini hesaplar. $T_{\text{base}}(y)$ sabit olduğundan state'te yer almaz.

    \item \textbf{Gauge fix (drift önleme):} Her zaman adımından sonra:
    \begin{equation}
        T' \leftarrow T' - \langle T' \rangle
        \label{eq:gauge_fix}
    \end{equation}
    Burada $\langle \cdot \rangle$ hacim ortalamasıdır. Neden? Periyodik domain'de ortalama sıcaklıkın zaman içinde kayması (numerical drift), sayısal hatalarda birikerek fiziksel olmayan sonuçlara yol açabilir. Bu düzeltme, $T'$'nin sıfır ortalamada kalmasını garanti eder.

    \item \textbf{Toplam sıcaklık:} Diagnostik amaçlı (Nusselt hesabı, görselleştirme):
    \begin{equation}
        T_{\text{total}}(x, y, z, t) = T_{\text{base}}(y) + T'(x, y, z, t).
    \end{equation}
\end{enumerate}

%\begin{figure}[H]
%    \centering
%    \includegraphics[width=0.90\textwidth]{temperature_decomposition.png}
%    \caption{Sıcaklık ayrıştırması. Sol: $T_{\text{base}}(y)$ lineer profil. Orta: $T'$ pertürbasyon (periyodik, FFT uyumlu). Sağ: $T_{\text{total}} = T_{\text{base}} + T'$.}
%    \label{fig:temperature_decomposition}
%\end{figure}


% ============================================================
\section{Boyutsuz Denklem Sistemi}
\label{sec:bizim_denklemler}

Proje, Boussinesq yaklaşımı altında çalışır (bkz.~Bölüm~\ref{ch:boussinesq}). Boyutsuz değişkenler ile denklem sistemi:

\subsection{Momentum Denklemi}
\begin{equation}
    \boxed{
        \pd{\vect{u}}{t} + (\vect{u} \cdot \gradient)\vect{u}
        = -\gradient p
        + \frac{1}{\Rey} \laplacian \vect{u}
        + \Rich \, T' \, \hat{\vect{e}}_y
        + \vect{F}
    }
    \label{eq:momentum_bizim}
\end{equation}

Terimlerin fiziksel anlamları:
\begin{itemize}
    \item $(\vect{u} \cdot \gradient)\vect{u}$: Adveksiyon --- akışkanın kendini taşıması (nonlineer terim).
    \item $-\gradient p$: Basınç gradyanı --- kütlenin korunumunu zorlar ($\divergence \vect{u} = 0$).
    \item $\frac{1}{\Rey} \laplacian \vect{u}$: Visköz difüzyon --- kinetik enerjiyi ısıya çevirir.
    \item $\Rich \, T' \, \hat{\vect{e}}_y$: Kaldırma kuvveti --- sıcaklık pertürbasyonundan kaynaklanan dikey kuvvet.
    \item $\vect{F}$: Dış zorlama (Kolmogorov forcing, (\ref{eq:kolmogorov_forcing})).
\end{itemize}

\subsection{Sıcaklık Pertürbasyon Denklemi}
\begin{equation}
    \boxed{
        \pd{T'}{t} + (\vect{u} \cdot \gradient)T'
        = \frac{1}{\Rey \, \Pran} \laplacian T'
        + v \, \frac{\Delta T}{L_y}
    }
    \label{eq:temperature_bizim}
\end{equation}
Son terim, Bölüm~\ref{sec:sicaklik_ayristirma}'daki türetimden gelen kaynak terimidir.

\subsection{Süreklilik Denklemi}
\begin{equation}
    \boxed{\divergence \vect{u} = 0}
    \label{eq:continuity_bizim}
\end{equation}
Boussinesq yaklaşımında hız alanı solenoidaldir. Bu koşul, her zaman adımında Projection nöronu tarafından sağlanır.

\subsection{Bileşen Formu}

Açık bileşen yazımı (hesaplama implementasyonu için):
\begin{align}
    \pd{u}{t} + u\pd{u}{x} + v\pd{u}{y} + w\pd{u}{z}
    &= -\pd{p}{x} + \frac{1}{\Rey}\left(\pdd{u}{x} + \pdd{u}{y} + \pdd{u}{z}\right) + F_x
    \label{eq:u_component} \\[6pt]
    \pd{v}{t} + u\pd{v}{x} + v\pd{v}{y} + w\pd{v}{z}
    &= -\pd{p}{y} + \frac{1}{\Rey}\left(\pdd{v}{x} + \pdd{v}{y} + \pdd{v}{z}\right) + \Rich \, T'
    \label{eq:v_component} \\[6pt]
    \pd{w}{t} + u\pd{w}{x} + v\pd{w}{y} + w\pd{w}{z}
    &= -\pd{p}{z} + \frac{1}{\Rey}\left(\pdd{w}{x} + \pdd{w}{y} + \pdd{w}{z}\right)
    \label{eq:w_component} \\[6pt]
    \pd{T'}{t} + u\pd{T'}{x} + v\pd{T'}{y} + w\pd{T'}{z}
    &= \frac{1}{\Rey \, \Pran}\left(\pdd{T'}{x} + \pdd{T'}{y} + \pdd{T'}{z}\right) + \frac{v \, \Delta T}{L_y}
    \label{eq:Tprime_component}
\end{align}
Önemli: Kaldırma kuvveti terimi \textit{yalnızca} $v$-denkleminde (dikey momentum) görülür; $u$ ve $w$ denklemlerinde yer almaz. Benzer şekilde, Kolmogorov forcing ($F_x$) yalnızca $u$-denklemindedir.


% ============================================================
\section{Boyutsuz Sayılar ve Parametre Uzayı}
\label{sec:parametre_uzayi}

\subsection{Boyutsuz Sayıların Tanımı}

\begin{table}[H]
\centering
\caption{Projede kullanılan boyutsuz sayılar.}
\label{tab:boyutsuz_sayilar}
\begin{tabular}{@{}llll@{}}
\toprule
\textbf{Sayı} & \textbf{Tanım} & \textbf{Değer Aralığı} & \textbf{Fiziksel Anlam} \\
\midrule
$\Rey$ & $UL/\nu$ & $5\,000$--$10\,000$ & Atalet / viskozite oranı \\
$\Ray$ & $g\beta\Delta T L^3 / (\nu\kappa)$ & $10^5$--$10^7$ & Kaldırma / difuzyon oranı \\
$\Pran$ & $\nu / \kappa$ & $0.71$ & Momentum / ısı difuzyon oranı (hava) \\
$\Rich$ & $\Ray / (\Rey^2 \Pran)$ & (hesaplanır) & Kaldırma / atalet oranı \\
$\Nus$ & $1 + \langle v T' \rangle \big/ \left(\kappa \frac{\Delta T}{L_y}\right)$ & $\geq 1$ & Toplam / difuzif ısı oranı \\
\bottomrule
\end{tabular}
\end{table}

Richardson sayısı akış rejimini belirleyen anahtar parametredir:
\begin{equation}
    \Rich = \frac{\Ray}{\Rey^2 \, \Pran}
    \label{eq:Ri_definition}
\end{equation}

\subsection{Richardson Sayısı Tablosu}

\begin{table}[H]
\centering
\caption{$\Rich$ değerleri, $\Rey$ ve $\Ray$ kombinasyonları için ($\Pran = 0.71$).}
\label{tab:Ri_values}
\begin{tabular}{@{}l|ccc@{}}
\toprule
$\Rey \backslash \Ray$ & $10^5$ & $10^6$ & $10^7$ \\
\midrule
$5\,000$   & $0.0056$   & $0.056$   & $0.56$     \\
$7\,000$   & $0.0029$   & $0.029$   & $0.29$     \\
$10\,000$  & $0.0014$   & $0.014$   & $0.14$     \\
\bottomrule
\end{tabular}
\end{table}

\subsection{Rejim Analizi}

\begin{itemize}
    \item $\Rich \ll 1$: \textbf{Rüzgâr baskın} (forced convection rejimi). Kolmogorov forcing, kaldırma kuvvetinden çok daha güçlü; sıcaklık pasif skaler gibi taşınır. Bu projede $\Rey = 10\,000$, $\Ray = 10^5$ ile $\Rich \approx 0.0014$: neredeyse tamamen zorlanmış konveksiyon.

    \item $\Rich \sim \mathcal{O}(0.1\text{--}1)$: \textbf{Mixed convection} (karışık rejim). Rüzgâr ve kaldırma kuvveti karşılaştırılabilir büyüklükte; en karmaşık fizik. Tilted plume'lar, eğik konveksiyon hücreleri görülür. Bu projede $\Rey = 5\,000$, $\Ray = 10^7$ ile $\Rich \approx 0.56$: mixed convection rejimi.

    \item $\Rich \gg 1$: \textbf{Kaldırma baskın} (natural convection rejimi). Sıcaklık farkı akışı domine eder; rüzgâr küçük bir pertürbasyondur. \textit{Bu projede seçilen parametre aralığında $\Rich > 1$ elde edilemiyor; en yüksek değer $\Rich \approx 0.56$.}
\end{itemize}

Tablodaki değerlere bakarsak, seçilen parametre uzayının $\Rich \sim 0.001$'den $\Rich \sim 0.56$'ya kadar uzandığını görürüz. Bu, zorlanmış konveksiyondan mixed konveksiyona geçişi kapsayan bir aralıktır. Kaldırma baskın rejim ($\Rich \gg 1$) için daha yüksek $\Ray$ veya daha düşük $\Rey$ değerleri gerekir.


% ============================================================
\section{CFL Koşulu ve Zaman Adımı}
\label{sec:cfl}

\subsection{CFL Koşulu}

Courant--Friedrichs--Lewy (CFL) koşulu, açık zaman integrasyon yöntemlerinin kararlılığı için gereklidir:
\begin{equation}
    \mathrm{CFL} = |\vect{u}|_{\max} \cdot \frac{\Delta t}{\Delta x_{\min}} < 1
    \label{eq:cfl}
\end{equation}
Üç boyutlu durumda:
\begin{equation}
    \mathrm{CFL} = \Delta t \left(\frac{|u|_{\max}}{\Delta x} + \frac{|v|_{\max}}{\Delta y} + \frac{|w|_{\max}}{\Delta z}\right) < 1.
    \label{eq:cfl_3d}
\end{equation}

Fiziksel anlam: bir zaman adımında bilgi, en fazla bir grid hücresi kadar ilerleyebilir. Bu koşul ihlal edilirse, sayısal şema dengesiz dalga modlarını üretir ve çözüm patlar.

\subsection{Grid Boyutları}

INNATE grid: $N_x \times N_y \times N_z = 96 \times 160 \times 64$.
\begin{equation}
    \Delta x = \frac{L_x}{N_x} = \frac{6}{96} = 0.0625, \quad
    \Delta y = \frac{L_y}{N_y} = \frac{10}{160} = 0.0625, \quad
    \Delta z = \frac{L_z}{N_z} = \frac{4}{64} = 0.0625.
    \label{eq:grid_spacing}
\end{equation}
Her üç yönde eşit grid aralığı ($\Delta x = \Delta y = \Delta z = 0.0625$) elde edilir. Bu izotropik grid, spektral yöntemler için idealdir.

DNS grid: $N_x \times N_y \times N_z = 192 \times 320 \times 128$. Grid aralığı: $\Delta x_{\text{DNS}} = 0.03125$.

\subsection{Zaman Adımı Analizi}

Seçilen zaman adımı: $\Delta t = 0.005$.

\textbf{Konvektif CFL:} Tipik türbülant hız $|\vect{u}|_{\max} \sim 3$ varsayımı ile:
\begin{equation}
    \mathrm{CFL}_{\text{conv}} = 3 \times \frac{0.005}{0.0625} = 0.24 \quad (\text{güvenli}).
    \label{eq:cfl_check}
\end{equation}

\textbf{Difüzif CFL:} 3B'de altı komşu nokta olduğundan:
\begin{equation}
    \Delta t < \frac{\Delta x^2}{6\nu} = \frac{0.0625^2}{6 / \Rey}
    \label{eq:diffusive_cfl}
\end{equation}
$\Rey = 5000$ ($\nu = 2 \times 10^{-4}$) için:
\[
    \Delta t < \frac{3.906 \times 10^{-3}}{1.2 \times 10^{-3}} \approx 3.26.
\]
Difüzif kararlılık sınırı çok rahat sağlanır ($\Delta t = 0.005 \ll 3.26$). Bu beklenendir: yüksek $\Rey$'de konvektif CFL kısıtlayıcıdır, difüzif değil.


% ============================================================
\section{Başlangıç Koşulları}
\label{sec:baslangic}

Simülasyon aşağıdaki başlangıç koşulları ile başlatılır:

\subsection{Hız Alanı}
\begin{equation}
    u(x, y, z, 0) = \epsilon_u \cdot \eta_u(x, y, z), \quad
    v(x, y, z, 0) = \epsilon_v \cdot \eta_v(x, y, z), \quad
    w(x, y, z, 0) = \epsilon_w \cdot \eta_w(x, y, z)
    \label{eq:ic_velocity}
\end{equation}
Burada $\eta_u, \eta_v, \eta_w \sim \mathcal{N}(0, 1)$ standart normal dağılımlı rastgele alanlar ve $\epsilon \sim 0.01$ küçük bir genlik faktörüdür.

\textbf{Neden küçük pertürbasyon?} Sıfır hız alanı ($\vect{u} = \vect{0}$) bir denge durumudur, ancak \textit{kararsız} bir dengedir (yüksek $\Rey$ ve $\Ray$ için). Küçük pertürbasyon, doğal instabiliteleri tetikler ve simülasyonun gerçekçi türbülant duruma geçmesini sağlar.

\begin{remark}
Başlangıç pertürbasyonu sonrası Projection uygulanarak $\divergence \vect{u} = 0$ sağlanır. Rastgele alan dolaylı olarak divergence-free değildir; Projection bu hatayı düzeltir.
\end{remark}

\subsection{Sıcaklık Alanı}
\begin{equation}
    T'(x, y, z, 0) = \epsilon_T \cdot \eta_T(x, y, z), \quad \epsilon_T \sim 0.01
    \label{eq:ic_temperature}
\end{equation}
Toplam sıcaklık: $T(x, y, z, 0) = T_{\text{base}}(y) + T'(x, y, z, 0)$.

\subsection{Basınç Alanı}
\begin{equation}
    p(x, y, z, 0) = 0.
    \label{eq:ic_pressure}
\end{equation}
Basınçın sıfırdan başlaması sorun oluşturmaz: Projection nöronu her adımda basıncı $\divergence \vect{u} = 0$'ı sağlayacak şekilde yeniden hesaplar. Başlangıç değeri ilk adımda düzeltilir.

\subsection{Spin-Up Dönemi}

Başlangıç koşullarından istatistiksel kararlı duruma ulaşmak için bir ``spin-up'' dönemi gereklidir. Bu dönemde:
\begin{enumerate}
    \item Kolmogorov forcing, ortalama akışı kurar.
    \item Sıcaklık farkı, konvektif hücreleri geliştirir.
    \item Enerji kaskadı oluşur: büyük ölçeklerden küçük ölçeklere enerji akar.
    \item \"{I}statistiksel denge kurulur: enerji üretimi $\approx$ dissipasyon.
\end{enumerate}
Spin-up süresi probleme bağlıdır; tipik olarak $t_{\text{spin-up}} \sim 10$--$50$ büyük ölçek zaman birimi ($L_y / U$).


% ============================================================
\section{Bölüm Özeti}
\label{sec:problem_ozet}

Bu bölümde tanımlanan problem:

\begin{itemize}
    \item $6 \times 10 \times 4$ boyutlu, üç yönde periyodik kutu.
    \item Alt sıcak ($20\,^{\circ}\mathrm{C}$), üst soğuk ($0\,^{\circ}\mathrm{C}$) $\to$ kaldırma kuvveti.
    \item Kolmogorov forcing $\to$ yatay rüzgâr etkisi.
    \item Sıcaklık ayrıştırması: $T = T_{\text{base}} + T'$ ile periyodik BC uyumu.
    \item 4 alan değişkeni: $u, v, w, T'$ (basınç Projection ile hesaplanır).
    \item 3 kontrol parametresi: $\Rey \in [5000, 10000]$, $\Ray \in [10^5, 10^7]$, $\Pran = 0.71$.
    \item Grid: $96 \times 160 \times 64$ (INNATE), $192 \times 320 \times 128$ (DNS).
    \item $\Delta t = 0.005$, $\mathrm{CFL} \approx 0.24$.
\end{itemize}

Sonraki bölümde (Bölüm~\ref{ch:innate_mimari}), bu denklemlerin her teriminin nasıl bir INNATE nöronuna dönüştüğünü ve ağ yapısını göreceğiz.
